\documentclass[letterpaper,10pt,conference]{ieeeconf}
\IEEEoverridecommandlockouts
\overrideIEEEmargins

\usepackage{amsmath,amssymb,amsthm,mathtools,bm}
\usepackage[hidelinks]{hyperref}
\usepackage{enumitem}
\usepackage{microtype}
\usepackage{cite}
\usepackage{xcolor} % Load the package
\allowdisplaybreaks
\emergencystretch=1.2em
\usepackage{graphicx}
\newtheorem{assumption}{Assumption}
\newtheorem{definition}{Definition}
\newtheorem{corollary}{Corollary}

\newtheorem{proposition}{Proposition}
\newtheorem{theorem}{Theorem}
\usepackage{comment}
\newtheorem{remark}{Remark}
\usepackage{algorithm}
\usepackage{algpseudocode}
\newcommand{\R}{\mathbb{R}}
\newcommand{\N}{\mathbb{N}}
\newcommand{\PP}{\mathbb{P}}
\newcommand{\Bm}{B_{\gamma}^{\tilde f}}
\newcommand{\Bt}{B_{\gamma}^{f}}
\newcommand{\fapp}{\widetilde{f}}
\newcommand{\safe}{\mathcal{L}}
\newcommand{\fail}{\mathcal{K}}
\newcommand{\xv}{x}
\newcommand{\uv}{u}
\newcommand{\dv}{d}
\newcommand{\Uv}{\mathcal{U}}
\newcommand{\Dv}{\mathcal{D}}
\newcommand{\Xiv}{\Xi}
\newcommand{\Xrl}{\mathcal{X}}
\newcommand{\norm}[1]{\left\lVert #1 \right\rVert}
\newcommand{\ip}[2]{\left\langle #1,#2 \right\rangle}
\newcommand{\pos}[1]{\left[#1\right]_+}

\title{Mondrian Conformal Calibration of Safety Certificates under Approximate Dynamics}
\author{}

\begin{document}
\maketitle

\begin{abstract}
Safety filters can be combined with reinforcement-learning controllers to enable safer deployment. In this paper, we propose a Mondrian conformal shielding framework for settings where the true dynamics are not fully known and safety certification must rely on an approximate model. The method uses a control barrier–value function (CBVF), which combines Hamilton–Jacobi reachability with barrier-function filtering, to construct a model-based safety filter. Because the CBVF is computed from an approximate dynamics model, the filter can become optimistic when the true system differs from the model. Using data collected from the system, our method calibrates how this model mismatch affects the safety certificate across different operating regions. Conformal calibration constructs safety buffers with a prescribed coverage level, while Mondrian calibration adapts these buffers to different operating regions. The resulting filter retains the standard CBVF-QP structure, reduces optimism caused by local model mismatch, and provides a probabilistic safety guarantee for the true system.
\\\textcolor{blue}{simulation}
\end{abstract}
\noindent\textbf{Keywords:} Safe reinforcement learning, safety filters,
control barrier--value functions, Hamilton--Jacobi reachability, conformal
prediction, Mondrian conformal prediction.
\section{Introduction}
% check plaigerisom
% citations in appropriate placess
%check abrivation
% inro, background, previous works, novelty , paper organization
% no repeating
% ***** feasability? ***
%%%%%%%%%%%%%%%%%%%%%%%%%%%%
%\textcolor{blue}{}

\section{Problem Formulation and Preliminaries}


This section formulates the model-based shielding problem considered in the
paper. We first introduce the reinforcement-learning policy, which provides the
nominal control input during deployment. Since the policy does not guarantee
safety along every realized trajectory, its input is passed through a safety
filter. We then introduce the true continuous-time dynamics and the approximate
model available for safety certification. The approximate model is used to
construct a CBVF certificate and the corresponding baseline CBVF-QP safety
filter. This baseline filter serves as the foundation for the data-driven calibration
developed in the following sections.
%%%%%%%%%%%%%%%%%%%
\subsection{Safe reinforcement learning and the nominal controller}

The reinforcement-learning policy is trained in an infinite-horizon discounted
constrained Markov decision process (CMDP)
\cite{altman1999cmdp,achiam2017cpo,tessler2019rcpo}, specified by the tuple
\begin{equation}
    (\Xrl,\Uv,P,g,c,\rho,\lambda,d_c),
    \label{eq:cmdp_tuple}
\end{equation}
where $\Xrl\subseteq\R^n$ and $\Uv\subseteq\R^{n_u}$ are the state and action
spaces, respectively, $P:\Xrl\times\Uv\to\mathcal{P}(\Xrl)$ is the transition
kernel, with $\mathcal{P}(\Xrl)$ denoting the set of probability measures on
$\Xrl$, $g:\Xrl\times\Uv\to\R$ is the reward function,
$c:\Xrl\times\Uv\to\R_{\ge0}$ is the safety-cost function,
$\rho\in\mathcal{P}(\Xrl)$ is the initial-state distribution,
$\lambda\in[0,1)$ is the discount factor, and $d_c\in\R_{\ge0}$ is the
prescribed safety-cost budget.

For a policy $\pi:\Xrl\to\mathcal{P}(\Uv)$, its expected discounted cumulative
safety cost is defined as
\begin{equation}
    J_c^\pi(\rho)
    :=
    \mathbb{E}_{\rho,\pi,P}
    \left[
        \sum_{k=0}^{\infty}\lambda^k c(x_k,u_k)
    \right],
    \label{eq:cmdp_cost}
\end{equation}
where the expectation is taken over trajectories initialized according to
$\rho$ and generated by the policy $\pi$ and transition kernel $P$. The policy
is required to satisfy
\begin{equation}
    J_c^\pi(\rho)\le d_c .
    \label{eq:cmdp_constraint}
\end{equation}
Although this constraint limits the expected cumulative safety cost, it does
not ensure pointwise safety along every realized trajectory. Therefore, a
deployment-time safety shield will later be used to modify the input proposed
by the learned policy when necessary
\cite{alshiekh2018safe,fisac2019general}.

At deployment, the learned policy provides the nominal controller
\begin{equation}
    \pi_{\mathrm{nom}}:\Xrl\times[t_0,0]\to\Uv,
    \label{eq:nominal_policy}
\end{equation}
where $[t_0,0]$ is the finite deployment horizon. For a stationary policy,
$\pi_{\mathrm{nom}}(\xv,t)=\pi_{\mathrm{nom}}(\xv)$. This nominal controller
will subsequently be used as the reference input for the safety filter
developed in the following sections.

%%%%%%%%%%%%%%%



\subsection{Continuous-time safety model}

We now introduce the continuous-time system and its safety specification.
Certification is carried out on a compact operating domain
$\Omega\subseteq\Xrl\subseteq\R^n$. The corresponding state--time domain over
the finite safety horizon is
\begin{equation}
    \Omega_T:=\Omega\times[t_0,0].
    \label{eq:state_time_domain}
\end{equation}
The true deployed dynamics are
\begin{equation}
    \dot{\xv}(s)=f(\xv(s),\uv(s),\dv(s)),
    \ \ \ \ 
    \xv\in\R^n,
    \uv\in\Uv,
    \dv\in\Dv\subset\R^{n_d},
    \label{eq:true_dynamics}
\end{equation}
while the safety certificate is computed with the approximate model
\begin{equation}
    \dot{\xv}(s)=\fapp(\xv(s),\uv(s),\dv(s)).
    \label{eq:model_dynamics}
\end{equation}

The true dynamics $f$ in \eqref{eq:true_dynamics} are not assumed to be known. Only $\fapp$ is available for
offline certification. This model mismatch is not included in the disturbance set $\Dv$, since this
would require a known bound on the mismatch and could make the resulting safety
certificate overly conservative.

Safety is specified by $l:\R^n\to\R$, with
\begin{equation}
    \safe:=\{\xv:l(\xv)\ge0\}.
    \label{eq:safe_set}
\end{equation}
The objective is to ensure that the true trajectory remains in $\safe$ over $[t_0,0]$ with a prescribed probability. We
impose the following standard regularity conditions on the dynamics and safety
specification.

\begin{assumption}[Standing regularity]
\label{ass:regularity}
The sets $\Uv$ and $\Dv$ in \eqref{eq:true_dynamics} are nonempty and compact.
The functions $f$ in \eqref{eq:true_dynamics} and $\fapp$ in
\eqref{eq:model_dynamics} are continuous in $(\xv,\uv,\dv)$ and locally
Lipschitz in $\xv$ on an open neighborhood of $\Omega$, uniformly over
$(\uv,\dv)\in\Uv\times\Dv$. The target function $l$ defining the safe set in
\eqref{eq:safe_set} is locally Lipschitz.
\end{assumption}

Assumption~\ref{ass:regularity} ensures local well-posedness of the
trajectories considered in the analysis. Under these conditions, we now introduce the safety certificate used to
construct the model-based filter.



%%%%


\subsection{Model-based control barrier--value function}

We now define the model-based CBVF, which is computed using the approximate
dynamics $\fapp$ in \eqref{eq:model_dynamics} and used for finite-horizon safety
certification \cite{choi2021cbvf}. For each $t\in[t_0,0]$, let $\Uv[t,0]$ and
$\Dv[t,0]$ denote the sets of measurable control and disturbance signals over
$[t,0]$, respectively. The disturbance is selected through a nonanticipative
strategy $\xi_{\dv}:\Uv[t,0]\to\Dv[t,0]$, and $\Xiv[t,0]$ denotes the
corresponding strategy set \cite{mitchell2005time}. For a CBVF discount rate $\gamma\in\R_{\ge0}$, define
\begin{equation}
\begin{aligned}
    \Bm(\xv,t)
    :=
    \inf_{\xi_{\dv}\in\Xiv[t,0]}
    \sup_{\uv(\cdot)\in\Uv[t,0]}
    \min_{\tau\in[t,0]}
    e^{\gamma(\tau-t)}l(\tilde{\xv}(\tau)),
\end{aligned}
\label{eq:model_cbvf}
\end{equation}
where $\tilde{\xv}(\cdot)$ denotes the trajectory of the approximate model
initialized at $\tilde{\xv}(t)=\xv$ and satisfying
\begin{equation}
    \dot{\tilde\xv}(\tau)=
    \fapp(\tilde\xv(\tau),\uv(\tau),\xi_{\dv}[\uv](\tau)).
    \label{eq:model_cbvf_traj}
\end{equation}
Under the standard Isaacs condition, $\Bm$ is characterized on the state--time
domain $\Omega_T$ defined in \eqref{eq:state_time_domain} as the viscosity
solution of
\begin{equation}
\label{eq:model_cbvf_vi}
\begin{aligned}
0=\min\Bigl\{\, &l(\xv)-\Bm(\xv,t),\;
\partial_t \Bm(\xv,t)+\gamma\Bm(\xv,t) \\
&+\max_{\uv\in\Uv}\min_{\dv\in\Dv}
\ip{\nabla_{\xv}\Bm(\xv,t)}{\fapp(\xv,\uv,\dv)}
\Bigr\},
\end{aligned}
\end{equation}
with terminal condition
\begin{equation}
    \Bm(\xv,0)=l(\xv).
    \label{eq:model_cbvf_terminal}
\end{equation}
Under Assumption~\ref{ass:regularity}, $\Bm$ is locally Lipschitz on
$\Omega_T$ and hence differentiable almost everywhere.
Since the minimum in \eqref{eq:model_cbvf} includes $\tau=t$,
\begin{equation}
    \Bm(\xv,t)\le l(\xv),
    \qquad (\xv,t)\in\Omega_T .
    \label{eq:Bm_le_l}
\end{equation}
Thus $\Bm(\xv,t)\ge0$ implies $\xv\in\safe$, and the model-based certified set is
\begin{equation}
    \mathcal C_{\gamma}^{\tilde f}(t)
    :=
    \{\xv\in\Omega:\Bm(\xv,t)\ge0\}.
    \label{eq:model_safe_set}
\end{equation}
The next subsection uses the CBVF certificate to construct the model-based
safety filter.



%%%%%%%
\subsection{Model-based safety filter}

We now use the model-based CBVF in \eqref{eq:model_cbvf} to construct an online
safety filter for the nominal controller in \eqref{eq:nominal_policy}
\cite{choi2021cbvf}. To formulate this filter as a convex quadratic program,
assume that the approximate model in \eqref{eq:model_dynamics} is
control-affine in the input on the operating domain and that $\Uv$ is compact
and convex. The uncalibrated model-based CBVF-QP is
\begin{equation}
\label{eq:model_cbvf_qp}
\begin{aligned}
\pi_{\mathrm{CBVF}}^{\tilde f}(\xv,t)
=\;& \arg\min_{\uv\in\Uv}
\norm{\uv-\pi_{\mathrm{nom}}(\xv,t)}^2 \\
\text{s.t.}\;&
\partial_t \Bm(\xv,t)+\gamma\Bm(\xv,t) \\
&+\min_{\dv\in\Dv}
\ip{\nabla_{\xv}\Bm(\xv,t)}{\fapp(\xv,\uv,\dv)}
\ge0 .
\end{aligned}
\end{equation}
The objective is strictly convex, and the feasible set is convex under the
control-affine structure. At differentiability points of $\Bm$, feasibility of
the baseline CBVF-QP on the model-based certified set in
\eqref{eq:model_safe_set} follows from the CBVF variational inequality
\eqref{eq:model_cbvf_vi}. Since the calibrated filter introduced later tightens
the CBVF constraint, its feasibility is considered separately.

At the decision times $t_0<t_1<\cdots<t_H=0$, the model-based safety filter
returns, for $j=0,\ldots,H-1$, the applied input
\begin{equation}
    u_j
    =
    \pi_{\mathrm{CBVF}}^{\tilde f}\!\left(\xv(t_j),t_j\right),
    \label{eq:shield_map}
\end{equation}
which is held constant over the corresponding sampling interval,
\begin{equation}
    \uv(s)=u_j,
    \qquad s\in[t_j,t_{j+1}).
    \label{eq:zoh_control}
\end{equation}
Since the filter is constructed using the approximate dynamics in
\eqref{eq:model_dynamics}, model mismatch may make its CBVF constraint too
optimistic for the true dynamics in \eqref{eq:true_dynamics}. The next section
explains how model mismatch affects the filter and how data-driven calibration
is used to account for it.



%%%%%%%%%%%%%%%%%%%%%%%%

\section{Safety Certification under Model Mismatch}
\label{sec:safety_model_mismatch}

This section identifies the component of model mismatch that can make the
model-based CBVF condition overly optimistic. The CBVF-QP filter in \eqref{eq:model_cbvf_qp} evaluates an input through the
CBVF rate. Since the time derivative $\partial_t\Bm$ is independent of the
system dynamics, model mismatch affects the CBVF condition through the
dynamics-dependent term
$\ip{\nabla_{\xv}\Bm(\xv,t)}{\dot\xv}$.
For a state-time point $(\xv,t)$, applied input $\uv$, and realized velocity
$\dot\xv$, we therefore define the directional model mismatch as
\begin{equation}
\begin{aligned}
    \mu(\xv,t,\uv,\dot\xv)
    :=
    &\min_{\dv\in\Dv}
    \ip{\nabla_{\xv}\Bm(\xv,t)}
        {\fapp(\xv,\uv,\dv)} \\
    &-
    \ip{\nabla_{\xv}\Bm(\xv,t)}
        {\dot\xv}.
\end{aligned}
\label{eq:directional_mismatch}
\end{equation}
The model-based quantity
$\min_{\dv\in\Dv}
\ip{\nabla_{\xv}\Bm(\xv,t)}{\fapp(\xv,\uv,\dv)}$
is computed from the approximate dynamics and the CBVF gradient by minimizing
over $\Dv$. The realized quantity
$\ip{\nabla_{\xv}\Bm(\xv,t)}{\dot\xv}$
uses measured or estimated system velocity. 
A positive value of $\mu$ in \eqref{eq:directional_mismatch} means that the approximate model predicts a more
favorable CBVF evolution than the true system, potentially making the
model-based constraint easier to satisfy than the true safety condition.
Values with $\mu\leq0$ do not weaken the safety condition. Moreover,
$\Bm(\xv,t)\ge0$ implies $l(\xv)\ge0$ by \eqref{eq:Bm_le_l}, independently of
the system dynamics. It is therefore sufficient to prevent the true
closed-loop trajectory from driving $\Bm$ below zero. The next section
introduces a data-driven method for accounting for positive directional
mismatch in the CBVF safety filter.




%%%%%%

\section{Mondrian Conformal Calibration of the CBVF Filter}
\label{sec:conformal_calibration}



This section develops a data-driven method for accounting for model mismatch in
the CBVF safety filter. We first use rollout data to measure how much the
approximate model overestimates the realized rate of change of the CBVF. We
then use these measurements to construct a safety buffer that tightens the
CBVF condition and compensates for this optimistic mismatch.

Because the mismatch may vary across the state--time domain, a single global
buffer may be unnecessarily conservative in some areas and insufficiently
responsive in others. We therefore partition the state--time domain into
regions and calibrate a separate buffer in each region. This region-wise
procedure is known as Mondrian conformal prediction
\cite{vovk2013conditional,bostrom2021mondrian}. Within each region, conformal
calibration uses the observed mismatch scores to determine a suitable buffer
without assuming a specific parametric distribution for the mismatch.


%%%%%%%%



\subsection{Pointwise mismatch score}
\label{subsec:pointwise_mismatch_score}

This subsection defines a one-sided score for measuring directional model
mismatch along closed-loop rollouts. For rollout $i$, let $\xv_i(s)$ denote
the state, $u_{i,j}$ the input held on $[t_j,t_{j+1})$, and
$\dot\xv_{i,j}(s)$ the corresponding measured or estimated state velocity.
The pointwise mismatch score is defined as
\begin{equation}
    \eta_{i,j}(s)
    :=
    \pos{\mu\bigl(\xv_i(s),s,u_{i,j},\dot\xv_{i,j}(s)\bigr)},
    \qquad s\in[t_j,t_{j+1}),
    \label{eq:eta_def}
\end{equation}
where $\mu$ is defined in \eqref{eq:directional_mismatch} and
$\pos{a}:=\max\{a,0\}$. The score retains only positive directional mismatch,
which can make the model-based CBVF condition optimistic.

If the approximate dynamics in \eqref{eq:model_dynamics} exactly match the
true dynamics in \eqref{eq:true_dynamics}, i.e., $f=\fapp$, and the realized
velocity is measured exactly, then $\eta_{i,j}(s)=0$. Otherwise, a positive
score captures the combined effect of model mismatch, unmodeled effects, and
velocity-estimation error on the CBVF condition. The next subsection groups
these pointwise scores by region and aggregates them over each rollout.


%%%%%%

\subsection{Regional partition and rollout scores}
\label{subsec:regional_scores}

This subsection partitions the operating domain and aggregates the pointwise
mismatch measurements into one score for each rollout and region. Let
$\{\Omega_m\}_{m=1}^M$ be a finite measurable partition of
$\Omega\subseteq\Xrl$ such that
\begin{equation}
    \Omega=\bigcup_{m=1}^M\Omega_m,
    \qquad
    \Omega_m\cap\Omega_{m'}=\emptyset
    \quad (m\ne m'),
    \label{eq:regional_partition}
\end{equation}
and let $m(\xv)\in\{1,\ldots,M\}$ denote the region containing $\xv$. The
partition may reflect obstacle geometry, terrain, or certificate properties
such as $\Bm$ and $\norm{\nabla_{\xv}\Bm}$.

For each decision interval, let
$\mathcal G_j\subset[t_j,t_{j+1})$ be the fixed evaluation-time grid used for
all calibration rollouts.

\begin{assumption}[Offline calibration dataset]
\label{ass:offline_calibration_rollouts}
The offline calibration dataset consists of $N_{\rm cal}$ rollouts
$\{(\xv_i,\uv_i,\dot\xv_i)\}_{i=1}^{N_{\rm cal}}$ over $[t_0,0]$. Each rollout
is generated by the true system in \eqref{eq:true_dynamics} using the
uncalibrated CBVF-QP in \eqref{eq:model_cbvf_qp} under an admissible
disturbance signal. For every $s\in\mathcal G_j$, $j=0,\ldots,H-1$, the
dataset contains the state $\xv_i(s)$, the applied input $u_{i,j}$, and the
corresponding measured or estimated state velocity $\dot\xv_{i,j}(s)$.
\end{assumption}

When the state velocity $\dot\xv$ is not directly measured, it can be estimated from
sampled state measurements using numerical differentiation or standard
state-estimation methods \cite{vanbreugel2020numerical,berkane2021nonlinear,
qin2018vinsmono}. The recorded quantities are used to compute the pointwise
score $\eta_{i,j}(s)$ in \eqref{eq:eta_def}.

A rollout $i$ visits region $m$ if $\xv_i(s)\in\Omega_m$ for some
$s\in\mathcal G_j$ and some $j$. The evaluation grids are chosen sufficiently
fine to capture every region visited by a rollout.

\begin{assumption}[Within-interval score regularity]
\label{ass:inter_sample_score}
Along every calibration rollout, the score $\eta_{i,j}$ in
\eqref{eq:eta_def} is Lipschitz continuous in time on each decision interval
with a common constant $L_\eta\in\mathbb{R}_{\geq 0}$. That is,
\begin{equation}
    |\eta_{i,j}(s)-\eta_{i,j}(s')|
    \leq L_\eta|s-s'|,
    \qquad
    s,s'\in[t_j,t_{j+1}).
    \label{eq:score_lipschitz}
\end{equation}
\end{assumption}

The constant $L_\eta$ may be derived from a uniform bound on
$|\dot\eta_{i,j}|$ over the decision intervals or conservatively estimated
from data \cite{wood1996estimation}.

Because the score is evaluated only on $\mathcal G_j$, its value may be larger
between evaluation times \cite{breeden2022sampled}. Let $h_{\rm eval}$ bound
the distance from any time at which a rollout lies in a region to the nearest
evaluation time in the same interval and region. By the choice of the evaluation grids $\mathcal G_j$, this bound can be taken as the largest grid spacing.
The possible inter-sample increase is then accounted for by choosing
\begin{equation}
    \epsilon_{\rm grid} = L_\eta h_{\rm eval},
    \label{eq:grid_margin}
\end{equation}
where $L_\eta$ is the within-interval Lipschitz constant in
\eqref{eq:score_lipschitz}, and $h_{\rm eval}\in\mathbb{R}_{\geq 0}$ is the
maximum time distance to the nearest evaluation point in the same interval
and region.


For each rollout $i$ that visits region $m$, define the regional rollout score
as
\begin{equation}
    z_{i,m}
    :=
    \max_{\substack{j\in\{0,\ldots,H-1\},\ s\in\mathcal G_j:\\
    \xv_i(s)\in\Omega_m}}
    \eta_{i,j}(s)
    +\epsilon_{\rm grid}.
    \label{eq:rollout_score}
\end{equation}
This score provides an upper bound on the continuous-time positive model mismatch
experienced by rollout $i$ in region $m$. For each region $m$, let
$\mathcal Z_m$ denote the set of regional rollout scores obtained from all
calibration rollouts that visit that region. Thus, each visiting rollout
contributes one score $z_{i,m}$ to $\mathcal Z_m$, and
$N_m:=|\mathcal Z_m|$ denotes the number of available calibration scores.
The corresponding regional score sets are used next to calibrate a conformal
model-mismatch buffer for each region.
%%%%%%%%%%%%%%%%%%

\subsection{Mondrian split-conformal calibration}

Here, we use Mondrian split-conformal calibration to construct a model-mismatch
buffer for each region. Let $\delta_{\rm traj}\in(0,1)$ denote the
trajectory-level miscoverage budget, which is the allowed miscoverage level for
the event that the model mismatch exceeds its calibrated buffer in at least one
region visited by a new rollout.

We distribute this budget across the $M$ regions by assigning each region a
miscoverage level $\delta_m\in(0,1)$ such that
\begin{equation}
    \sum_{m=1}^{M}\delta_m\leq\delta_{\rm traj}.
    \label{eq:regional_risk_allocation}
\end{equation}
We use the uniform allocation $\delta_m=\delta_{\rm traj}/M$, which assigns the
same miscoverage budget to each region. By
\eqref{eq:regional_risk_allocation}, the regional miscoverage probabilities can
be combined using a union bound, without assuming that buffer exceedances in
different regions occur independently
\cite{bostrom2021mondrian,dunn1961multiple}.
\begin{assumption}[Regional exchangeability]
\label{ass:offline_regional_exchangeability}
For each region $m$, among rollouts that visit that region, the calibration
scores in $\mathcal Z_m$ and the regional score of a new rollout are
exchangeable. The calibration and new rollouts are generated using the same
uncalibrated CBVF-QP in \eqref{eq:model_cbvf_qp}.
\end{assumption}

The exchangeability condition in
Assumption~\ref{ass:offline_regional_exchangeability} is imposed within each
region and does not require exchangeability across different regions.

We now use the regional calibration scores to compute the conformal
model-mismatch buffers. Let
$z_{(1),m}\leq\cdots\leq z_{(N_m),m}$ denote the ordered regional scores from
\eqref{eq:rollout_score}, where $N_m$ is the number of calibration rollouts
that visited region $m$. For the regional miscoverage level $\delta_m$, define
\begin{subequations}
\begin{align}
    k_m
    &:=
    \left\lceil (N_m+1)(1-\delta_m)\right\rceil,
    \label{eq:offline_mondrian_quantile_index}
    \\
    \widehat\xi^{\rm off}_{m}
    &:=
    z_{(k_m),m}.
    \label{eq:offline_mondrian_buffer}
\end{align}
\end{subequations}
where $k_m\in\mathbb N$ is the index of the conformal quantile for region $m$,
and $\widehat\xi^{\rm off}_{m}\in\mathbb R_{\geq0}$ is the resulting calibrated
model-mismatch buffer.

This finite-sample construction requires
\begin{equation}
    N_m\geq
    \left\lceil\frac{1}{\delta_m}\right\rceil-1,
    \label{eq:regional_sample_requirement}
\end{equation}
which is equivalent to $\delta_m\geq1/(N_m+1)$
\cite{lei2018distributionfree}. Under the uniform allocation
$\delta_m=\delta_{\rm traj}/M$,
\eqref{eq:regional_sample_requirement} becomes
$N_m\geq\lceil M/\delta_{\rm traj}\rceil-1$. If this condition is not satisfied,
the partition is coarsened according to a predetermined rule, for example
based on spatial proximity or similar operating conditions, and calibration is
repeated using a fresh dataset on the resulting fixed partition.

We next establish the finite-sample regional coverage guarantee provided by the
calibrated buffers.

\begin{proposition}[Mondrian split-conformal calibration]
\label{prop:mondrian_initialization}
For each region $m$, let $V_m$ denote the event that a new rollout visits
region $m$, and let $Z_{{\rm new},m}$ denote its regional score when $V_m$
occurs. Define the set of regions visited by the rollout as
\begin{equation}
    \mathcal V:=
    \{m\in\{1,\ldots,M\}:V_m\text{ occurs}\}.
    \label{eq:visited_region_set}
\end{equation}
Suppose Assumption~\ref{ass:offline_regional_exchangeability} holds, the
sample-size condition in \eqref{eq:regional_sample_requirement} is satisfied
for every region, and the regional miscoverage levels satisfy
\eqref{eq:regional_risk_allocation}. Then, for every region $m$ with
$\PP(V_m)>0$,
\begin{equation}
    \PP\!\left(
        Z_{{\rm new},m}>\widehat\xi^{\rm off}_{m}
        \,\middle|\,V_m
    \right)
    \leq\delta_m,
    \label{eq:conditional_mondrian_coverage}
\end{equation}
and
\begin{equation}
    \PP\!\left(
        Z_{{\rm new},m}\leq\widehat\xi^{\rm off}_{m},
        \ \forall m\in\mathcal V
    \right)
    \geq1-\delta_{\rm traj}.
    \label{eq:trajectory_mondrian_coverage}
\end{equation}
\end{proposition}

\begin{proof}
For each region $m$, conditional exchangeability and the standard
split-conformal rank argument
\cite{vovk2013conditional,lei2018distributionfree} give
\begin{equation*}
    \PP\!\left(
        Z_{{\rm new},m}>\widehat\xi^{\rm off}_{m}
        \,\middle|\,V_m
    \right)
    \leq
    1-\frac{k_m}{N_m+1}
    \leq\delta_m.
\end{equation*}
Define the regional buffer-exceedance event
\[
    F_m:=
    V_m\cap
    \left\{
        Z_{{\rm new},m}>\widehat\xi^{\rm off}_{m}
    \right\}.
\]
Then
\[
    \PP(F_m)
    =
    \PP\!\left(
        Z_{{\rm new},m}>\widehat\xi^{\rm off}_{m}
        \,\middle|\,V_m
    \right)\PP(V_m)
    \leq\delta_m.
\]
The union bound together with \eqref{eq:regional_risk_allocation} therefore
gives
\begin{equation*}
    \PP\!\left(\bigcup_{m=1}^{M}F_m\right)
    \leq
    \sum_{m=1}^{M}\delta_m
    \leq\delta_{\rm traj}.
\end{equation*}
Taking complements gives \eqref{eq:trajectory_mondrian_coverage}.
\hfill\qed
\end{proof}

Proposition~\ref{prop:mondrian_initialization} establishes the finite-sample
coverage guarantee for the regional conformal buffers. These buffers are used
next to construct the calibrated CBVF safety filter.



%%%%%%%%%%%%%%%%%


\section{Calibrated CBVF Filter and Safety Analysis}
\label{sec:calibrated_filter_safety}

This section incorporates the regional conformal buffers, which quantify the
positive model mismatch observed in each operating region, into the CBVF safety
filter. We first use the buffers in \eqref{eq:offline_mondrian_buffer} to
calibrate the model-based CBVF-QP in \eqref{eq:model_cbvf_qp}. We then
describe the final conformal certification of the resulting calibrated filter
and establish the corresponding safety guarantee under the true dynamics.

%%%%%%%%%%%%%%%%%%
\subsection{Region-wise calibrated CBVF-QP}

At each decision time $t_j$, the filter in \eqref{eq:shield_map} computes an
input $u_j\in\Uv$, which is held constant over $[t_j,t_{j+1})$ according to
\eqref{eq:zoh_control}. Define
\begin{equation}
\begin{aligned}
\Psi(\xv,t,\uv)
:={}&
\partial_t \Bm(\xv,t)+\gamma\Bm(\xv,t) \\
&+
\min_{\dv\in\Dv}
\ip{\nabla_{\xv}\Bm(\xv,t)}
{\fapp(\xv,\uv,\dv)} .
\end{aligned}
\label{eq:psi_definition}
\end{equation}
The uncalibrated CBVF-QP in \eqref{eq:model_cbvf_qp} enforces
$\Psi(\xv,t,\uv)\geq0$. The calibrated filter instead incorporates the Mondrian conformal buffer in \eqref{eq:offline_mondrian_buffer} to compensate for regional model mismatch.

To account for possible state transitions between control updates in
\eqref{eq:shield_map}--\eqref{eq:zoh_control}, we impose the following
assumption that the inter-sample state variation is bounded.
\begin{assumption}[Inter-sample state bound]
\label{ass:bounded_inter_sample_motion}
There exists a known constant $C_x\in\mathbb R_{\geq0}$ such that the true
closed-loop trajectory satisfies
\begin{equation}
\norm{\xv(s)-\xv(t_j)}
\leq
C_x(s-t_j),
\qquad
s\in[t_j,t_{j+1}).
\label{eq:bounded_inter_sample_motion}
\end{equation}
\end{assumption}

The constant $C_x$ in \eqref{eq:bounded_inter_sample_motion} may be obtained
from available physical, actuator, hardware, or certified operating-envelope
bounds, without requiring knowledge of the exact true dynamics. Such
inter-sample bounds are standard in sampled-data safety analysis
\cite{breeden2022sampled}. Using the bound $C_x$ in \eqref{eq:bounded_inter_sample_motion}, we
identify the Mondrian regions that the state may enter before the next control
update. Specifically, define
\begin{equation}
\begin{aligned}
\mathcal M_j
:=
\biggl\{
m\in\{1,\ldots,M\}:\
&\inf_{\xv'\in\Omega_m}
\norm{\xv(t_j)-\xv'}_2 \\
&\leq C_x(t_{j+1}-t_j)
\biggr\}.
\end{aligned}
\label{eq:reachable_region_set}
\end{equation}
Thus, $\mathcal M_j$ contains all regions that are reachable from $\xv(t_j)$ within the inter-sample state bound. The partition and calibration dataset are
chosen such that every region that may belong to $\mathcal M_j$ for a
certified decision-time configuration satisfies the regional sample-size
requirement in \eqref{eq:regional_sample_requirement}.

Since the state may transition between these regions while the input is held
constant, the filter uses the largest corresponding regional buffer,
\begin{equation}
\widehat\xi^{\rm app}_{j}
:=
\max_{m\in\mathcal M_j}
\widehat\xi^{\rm off}_{m},
\label{eq:applied_reachable_buffer}
\end{equation}
which ensures that the applied buffer covers every Mondrian region that may be
entered during $[t_j,t_{j+1})$.

To account for the variation of the CBVF constraint between decision times, we
impose the following regularity condition.
\begin{assumption}[Inter-sample regularity]
\label{ass:inter_sample_psi}
The quantity $\Psi$ in \eqref{eq:psi_definition} is well defined on
$\Omega_T$. For each decision interval, a known constant
$L_{\Psi,j}\in\mathbb R_{\geq0}$ is available such that, uniformly over
$\uv\in\Uv$,
\begin{equation}
\left|
\Psi(\xv,t,\uv)-\Psi(\xv',t',\uv)
\right|
\leq
L_{\Psi,j}
\left(
\norm{\xv-\xv'}+|t-t'|
\right)
\label{eq:psi_lipschitz}
\end{equation}
for all state--time points relevant to the interval.
\end{assumption}

The constants $L_{\Psi,j}$ in \eqref{eq:psi_lipschitz} may be obtained
analytically from bounds on the state and time derivatives of $\Psi$, or
conservatively certified over the relevant inter-sample reachable set using
numerical or interval-based bounding methods
\cite{breeden2022sampled,singletary2020sampled,zhang2022interval}.

From Assumption~\ref{ass:inter_sample_psi} and
\eqref{eq:bounded_inter_sample_motion}, for
$s\in[t_j,t_{j+1})$,
\begin{equation}
\left|
\Psi(\xv(s),s,\uv)-\Psi(\xv(t_j),t_j,\uv)
\right|
\leq
L_{\Psi,j}(C_x+1)(t_{j+1}-t_j).
\label{eq:psi_inter_sample_bound}
\end{equation}
We therefore define the inter-sample margin
\begin{equation}
\epsilon_{{\rm inter},j}
:=
L_{\Psi,j}(C_x+1) (t_{j+1}-t_j) .
\label{eq:inter_sample_margin}
\end{equation}

The calibrated CBVF filter is
\begin{equation}
\begin{aligned}
u_j=\arg\min_{\uv\in\Uv}\;&
\norm{\uv-\pi_{\mathrm{nom}}(\xv(t_j),t_j)}^2 \\
\mathrm{s.t.}\;&
\Psi(\xv(t_j),t_j,\uv)
\geq
\widehat\xi^{\rm app}_{j}
+\epsilon_{{\rm inter},j}.
\end{aligned}
\label{eq:regional_cp_filter}
\end{equation}
Thus, the nominal CBVF-QP objective is retained, while its constraint is tightened by the regional model-mismatch buffer and the inter-sample margin. This tightening changes the closed-loop behavior relative to the uncalibrated filter used to construct the regional buffers in \eqref{eq:offline_mondrian_buffer}. Therefore, an additional certification step is required to verify that these buffers remain valid under the calibrated filter \eqref{eq:regional_cp_filter}.


\begin{remark}[Final conformal certification]
\label{rem:final_cp_certification}
The regional buffers in \eqref{eq:offline_mondrian_buffer} are obtained from
rollouts of the uncalibrated filter \eqref{eq:model_cbvf_qp}. Since using these
buffers in \eqref{eq:regional_cp_filter} changes the closed-loop policy,
exchangeability with the original calibration rollouts may no longer hold.

Let
\begin{equation}
\widehat\xi^{(0)}_{m}
:=
\widehat\xi^{\rm off}_{m},
\qquad
m=1,\ldots,M.
\end{equation}
At iteration $r\in\mathbb N_0$, collect fresh rollouts under the calibrated
filter using $\widehat\xi^{(r)}_{m}$ and apply the same regional conformal
construction in
\eqref{eq:offline_mondrian_quantile_index}--\eqref{eq:offline_mondrian_buffer}
to obtain $\widehat\xi^{(r+1)}_{m}$. The buffers are updated iteratively with
fresh rollouts until the newly computed regional quantiles do not exceed the
currently applied buffers. The resulting values $\widehat\xi^{(r)}_{m}$ are
then taken as the final $\widehat\xi^{\rm off}_{m}$ in
\eqref{eq:applied_reachable_buffer} and fixed for deployment. The final
certification rollouts and a new deployment rollout are assumed to be
exchangeable under the same calibrated filter.
\end{remark}







We impose the following feasibility condition on the tightened QP.

\begin{assumption}[Calibrated-QP feasibility]
\label{ass:calibrated_qp_feasibility}
At every decision time $t_j$, the calibrated QP in
\eqref{eq:regional_cp_filter} is feasible, i.e., there exists
$\uv\in\Uv$ satisfying its constraint.
\end{assumption}

When this condition is not satisfied, feasibility can instead be addressed
using standard modifications such as backup-controller constructions or
adaptive relaxation of the barrier condition
\cite{chen2021backup,zeng2021optimal}. Incorporating such mechanisms into the
proposed conformal calibration framework is beyond the scope of this paper.

Now that the calibrated filter has been constructed, we next establish the corresponding closed-loop safety guarantee.

















%%%%%%%%%%%%%%%%%%%%


\subsection{Safety analysis}
Here, we establish the main probabilistic safety guarantee for the RL closed-loop system under the calibrated CBVF filter using conformal certification.


\begin{theorem}[Probabilistic closed-loop safety under model mismatch]
\label{thm:covered_mismatch_safety}
%%%%%%%%%%%%%%%%%%%%%
Consider the continuous-time system $f$ in \eqref{eq:true_dynamics}, operated by
the nominal RL controller $\pi_{\mathrm{nom}}$ in
\eqref{eq:nominal_policy} through the calibrated CBVF safety filter
\eqref{eq:regional_cp_filter}. The filter is constructed using the
approximate model $\fapp$ in \eqref{eq:model_dynamics}. Suppose the initial
condition lies in the model-based certified set, i.e.,
$\xv(t_0)\in\mathcal C_{\gamma}^{\tilde f}(t_0)$ defined in
\eqref{eq:model_safe_set}.
%%%%%%%%%%
Then the probability that the closed-loop state trajectory under the true
dynamics in \eqref{eq:true_dynamics} remains in the safe set over $[t_0,0]$
satisfies
\begin{equation}
\PP\!\left(
\xv(s)\in\safe,\ \forall s\in[t_0,0]
\right)
\geq
1-\delta_{\rm traj}.
\label{eq:calibrated_probabilistic_safety}
\end{equation}
\end{theorem}
%%%%%%%%%%%%%%
\begin{proof}
By the final conformal certification in
Remark~\ref{rem:final_cp_certification}, together with the regional risk
allocation in \eqref{eq:regional_risk_allocation}, with probability at least
$1-\delta_{\rm traj}$,
\[
Z_{{\rm new},m}
\leq
\widehat\xi^{\rm off}_{m}
\]
for every region $m$ visited by the deployment rollout.
%%%%%%%%%%%
Consider a deployment rollout for which this coverage event holds. For
$s\in[t_j,t_{j+1})$, let $m_s:=m(\xv(s))$. Since
$m_s\in\mathcal M_j$ by \eqref{eq:reachable_region_set}, for almost every
$s\in[t_j,t_{j+1})$,
\[
\eta(s)
\overset{\eqref{eq:rollout_score}}{\leq}
Z_{{\rm new},m_s}
\overset{\eqref{eq:trajectory_mondrian_coverage}}{\leq}
\widehat\xi^{\rm off}_{m_s}
\overset{\eqref{eq:applied_reachable_buffer}}{\leq}
\widehat\xi^{\rm app}_{j}.
\tag{*}
\]

The calibrated filter \eqref{eq:regional_cp_filter} enforces the tightened
CBVF constraint at the sampling time $t_j$,
\[
\Psi(\xv(t_j),t_j,u_j)
\geq
\widehat\xi^{\rm app}_{j}
+
\epsilon_{{\rm inter},j}.
\]
The additional margin $\epsilon_{{\rm inter},j}$ accounts for the possible
variation of $\Psi$ between consecutive sampling times. Hence, by
\eqref{eq:psi_inter_sample_bound} and \eqref{eq:inter_sample_margin},
\[
\Psi(\xv(s),s,u_j)
\geq
\widehat\xi^{\rm app}_{j},
\qquad
s\in[t_j,t_{j+1}).
\]
Combining this with $(*)$ gives
\[
\Psi(\xv(s),s,u_j)\geq\eta(s)
\tag{**}
\]
for almost every $s\in[t_j,t_{j+1})$.
%%%%%%%%%%%%
Since $\eta(s)=[\mu(s)]_+$, we have $\mu(s)\leq\eta(s)$. Using
\eqref{eq:directional_mismatch} and \eqref{eq:psi_definition}, it follows
that, for almost every $s\in[t_j,t_{j+1})$,
\[
\frac{d}{ds}\Bm(\xv(s),s)
+\gamma\Bm(\xv(s),s)
\geq
\Psi(\xv(s),s,u_j)-\eta(s)
\overset{(**)}{\geq}0.
\tag{***}
\]
%%%%%%%%%%%%
\textcolor{red}{here}
From (***) and
Assumption~\ref{ass:regularity},
\[
e^{\gamma(s-t_0)}\Bm(\xv(s),s)
\]
is nondecreasing on $[t_0,0]$. Since
$\xv(t_0)\in\mathcal C_{\gamma}^{\tilde f}(t_0)$,
\eqref{eq:model_safe_set} gives
$\Bm(\xv(t_0),t_0)\geq0$, and therefore
\[
\Bm(\xv(s),s)\geq0,
\qquad
s\in[t_0,0].
\]
By \eqref{eq:Bm_le_l}, this implies
$l(\xv(s))\geq0$, and hence
$\xv(s)\in\safe$ for all $s\in[t_0,0]$.
%%%%%%%%%%%%%
Since the final conformal certification holds with probability at least
$1-\delta_{\rm traj}$, \eqref{eq:calibrated_probabilistic_safety} follows.
%%%%%%%%%%%%%%%
\hfill\qed
\end{proof}






%%%%%%%%%%%%%%%









%%%%%%%%%%%%%%%%%%%%%

\section{Implementation and Experimental Validation}
\label{sec:experiments}
\begin{comment}
% =====================================================================
%  Easy-mismatch experiment: Mondrian conformal calibration of a
%  CBVF safety filter for shielded RL (Dubins car reach--avoid).
%
%  Self-contained section. Requires: booktabs, amsmath, siunitx (optional).
%  If you don't use siunitx, the numbers are plain and will still compile.
% =====================================================================

\section{Experiment: Reach--Avoid under Approximate Dynamics (Easy Mismatch)}
\label{sec:easy-mismatch}

\subsection{Setup}

We evaluate the proposed Mondrian split-conformal safety filter on a
Dubins-car reach--avoid task. The agent is trained with
SAC-Lagrangian~\cite{ha2021} under a mandatory safety shield derived from a
control barrier--value function (CBVF)~\cite{choi2021} that is itself computed
offline from an \emph{approximate} model of the dynamics. The safety
certificate is tightened online by region-wise conformal buffers
$\hat\xi$ calibrated from observed model--truth mismatch, so that the
executed policy remains safe despite the model error.

The Dubins car has state $(x,y,\theta)$ and turn-rate control
$u\in[-1,1]$. The safety constraint enforced by the shield at each active
state is
\begin{equation}
\label{eq:cbvf-constraint}
\partial_t B(x)
\;+\; \langle \nabla B(x),\, f_{\mathrm{drift}}(x)\rangle
\;+\; \gamma\, B(x)
\;+\; a_u(x)\,u
\;\ge\; \hat\xi_{r(x)},
\end{equation}
where $B$ is the CBVF value, $\gamma$ the CBVF discount, $a_u=\langle
\nabla_\theta B,\beta\rangle$ the control coefficient, and
$\hat\xi_{r(x)}$ the conformal buffer for the Mondrian region $r(x)$
containing the state. The shield activates only inside the certified set
and near its boundary; where it activates it solves a one-dimensional
quadratic program that projects the nominal action onto the half-space
defined by~\eqref{eq:cbvf-constraint}.

\paragraph{Model--truth mismatch (easy).}
The CBVF is computed for a model with turn-authority
$\beta_{\mathrm{model}}=0.75$ while the true system has
$\beta_{\mathrm{truth}}=0.60$; both share speed $v=0.60$. The filter
therefore certifies safety using a control-effectiveness estimate that
is optimistic relative to the truth, which is exactly the regime the
conformal calibration is designed to correct.

\paragraph{Mondrian calibration.}
The state space is partitioned by signed obstacle clearance into four
regions with edges $\{0.25,0.75,1.50\}$. Each region receives its own
split-conformal buffer at miscoverage level $\delta=0.05$, calibrated
from $n_{\mathrm{calib}}=500$ shielded rollouts of the baseline
(uncalibrated) CBVF filter.

\begin{table}[t]
\centering
\caption{Key configuration for the easy-mismatch experiment. All three
seeds use identical settings.}
\label{tab:easy-config}
\begin{tabular}{@{}ll@{}}
\toprule
Parameter & Value \\
\midrule
Training algorithm         & SAC-Lagrangian \\
Training steps             & $5\times10^{5}$ \\
Evaluation episodes        & 50 (per method, per eval) \\
CBVF grid                  & $81\times81\times61$ \\
CBVF time step             & $0.10$ \\
CBVF discount $\gamma$     & $0.10$ \\
Integration step $dt$      & $0.05$, horizon $400$ \\
Model $(v,\beta)$          & $(0.60,\,0.75)$ \\
Truth $(v,\beta)$          & $(0.60,\,0.60)$ \\
Mondrian clearance edges   & $\{0.25,0.75,1.50\}$ \\
Calibration rollouts $n_{\mathrm{calib}}$ & $500$ \\
Miscoverage $\delta$       & $0.05$ \\
Filter activation margin   & $0.40$ \\
Regional buffer ceiling    & $0.10$ \\
\bottomrule
\end{tabular}
\end{table}

\subsection{Design choices for feasibility and robustness}
\label{sec:design-choices}

Three design decisions ensure the calibrated quadratic program is
feasible where the policy actually operates, and that feasibility does
not depend on the random calibration draw.

\paragraph{Activation depth.}
The shield activates at a margin of $0.40$ inside the certified set
rather than at its boundary $B\approx 0$. At the certified-set boundary
the certificate is steep in time (the reach--avoid value changes rapidly
near the terminal horizon), so the control authority available there is
insufficient to meet the required buffer and the projection QP becomes
infeasible. Activating deeper in the safe set engages the filter where
the vehicle still has the turning authority to satisfy
\eqref{eq:cbvf-constraint}, which is also the correct place to intervene:
$B\approx 0$ is the last-ditch edge of the guarantee, not the obstacle
surface. This is a deliberate design choice, not a tuned nuisance
parameter.

\paragraph{Bounded-rate feasibility gate.}
The offline feasibility audit sweeps every certified grid node. A small
set of near-terminal states (largest at CBVF time $\approx -0.2$) exhibit
a finite-difference $\partial_t B$ spike and cannot meet the buffer;
these states are rarely if ever occupied by the trained policy at
runtime. We therefore gate on a \emph{bounded} offline infeasibility
rate ($\le 40\%$ of active certified nodes) rather than requiring zero
infeasible nodes, consistent with the exceedance semantics of conformal
prediction. The operationally meaningful quantity is the \emph{runtime}
infeasibility rate, which we report and which is small (see below).

\paragraph{Buffer ceiling.}
Far from the obstacle, model--truth trajectories diverge over the long
horizon and produce large mismatch scores. Because the filter never
needs to act in that region, an uncapped conformal quantile there can
inflate the buffer and, with an unlucky calibration draw, render the
audit infeasible. We cap each regional buffer at $0.10$. The
near-obstacle regions --- where safety is actually enforced --- calibrate
stably at $\hat\xi\approx 0.058$ and $0.063$ across seeds and are
unaffected by the cap; only the far-field region is clipped.

\subsection{Results}
\label{sec:easy-results}

We compare three controllers under the true dynamics: the unshielded
nominal policy (\textsc{nominal}), the policy with the uncalibrated CBVF
filter (\textsc{cbvf}), and the policy with the proposed conformally
calibrated filter (\textsc{cp}). Table~\ref{tab:easy-main} reports the
final-evaluation metrics over 50 episodes for a representative seed;
Table~\ref{tab:easy-seeds} reports the \textsc{cp} controller across
three seeds.

\begin{table}[t]
\centering
\caption{Final evaluation under the true dynamics (easy mismatch, seed~1,
50 episodes). \textsc{cp} is the only method that is both collision-free
and goal-reaching.}
\label{tab:easy-main}
\begin{tabular}{@{}lccc@{}}
\toprule
Method & Unsafe rate $\downarrow$ & Goal rate $\uparrow$ & Cost $\downarrow$ \\
\midrule
\textsc{nominal} & 0.92 & 0.00 & 0.16 \\
\textsc{cbvf}    & 0.92 & 0.00 & 0.18 \\
\textsc{cp} (ours) & \textbf{0.00} & \textbf{1.00} & \textbf{0.00} \\
\bottomrule
\end{tabular}
\end{table}

\begin{table}[t]
\centering
\caption{Seed stability of the proposed \textsc{cp} controller (easy
mismatch, 50 episodes each). Safety is invariant across seeds; goal rate
and runtime infeasibility vary with the calibration draw, and the method
degrades gracefully --- an unluckier draw (seed~3) loses some goal
completions but never safety.}
\label{tab:easy-seeds}
\begin{tabular}{@{}lccc@{}}
\toprule
Seed & Unsafe rate & Goal rate & Runtime infeas./active \\
\midrule
1 & 0.00 & 1.00 & 0.002 \\
2 & 0.00 & 1.00 & 0.001 \\
3 & 0.00 & 0.76 & 0.048 \\
\midrule
Range & $0.00$ & $0.76$--$1.00$ & $0.001$--$0.048$ \\
\bottomrule
\end{tabular}
\end{table}

\paragraph{Safety.}
Across all three seeds the calibrated filter yields a
$0.00$ unsafe rate, i.e.\ zero collisions in $150$ evaluation episodes,
whereas both the nominal policy and the uncalibrated CBVF filter collide
in $92$--$100\%$ of episodes. The conformal calibration is what converts
an optimistic, uncalibrated certificate that fails almost always into one
that is safe under the true dynamics.

\paragraph{Feasibility.}
The calibrated projection QP is feasible at essentially every state the
policy visits: the runtime infeasibility rate among active steps is
$0.1$--$0.2\%$ on seeds~1 and~2 and $4.8\%$ on seed~3. On the rare
infeasible step the filter falls back to a least-violation actuator-rail
action, which is logged. On feasible active steps the filter performs the
exact minimum-distance projection
(\texttt{minimal\_projection\_verified}), so the reported safety is
produced by the genuine calibrated QP rather than by the fallback.

\paragraph{Task performance.}
The calibrated filter reaches the goal in every episode on seeds~1 and~2
and in $76\%$ of episodes on seed~3, at zero safety cost. The lower goal
rate on seed~3 is a direct consequence of its higher runtime
infeasibility: additional rail fallbacks perturb the trajectory and cost
some goal completions, but never induce a collision. The nominal and
uncalibrated CBVF controllers reach the goal in $0\%$ of episodes under
this mismatch.

\subsection{Summary}

Under an optimistic model--truth mismatch on the Dubins-car reach--avoid
task, the proposed Mondrian conformal filter is the only evaluated method
that is simultaneously collision-free ($0\%$ unsafe across three seeds)
and goal-reaching (goal rate $0.76$--$1.00$), with runtime QP
infeasibility of at most $4.8\%$. Safety is invariant to the calibration
seed; the remaining variability is confined to task performance and
degrades gracefully.
\end{comment}

\section{Conclusion}

This paper introduced a Mondrian conformal calibration framework for
CBVF-based safety filtering under approximate dynamics. We used a model-based
CBVF certificate as the baseline safety mechanism and calibrated its
derivative condition using rollout data. The calibration was performed region
wise, allowing the safety buffer to reflect local variation in model mismatch.

We used Mondrian split conformal calibration to construct a data-driven
region-wise buffer while preserving the standard CBVF-QP structure. The analysis separated the statistical calibration layer from the deterministic
CBVF safety argument. Final conformal certification is performed under the
calibrated filter so that the certification and deployment rollouts are
exchangeable under the same closed-loop policy, yielding the resulting
deployment-time probabilistic safety guarantee. \\\textcolor{blue}{simulation}

\bibliographystyle{IEEEtran}
\bibliography{conformal_cbvf_shielding}

\end{document}
